\documentclass[12pt]{article}

% Use algorithm and algpseudocode packages so \State and related commands are available

\newcommand{\duedate}{---}
\newcommand{\assignment}{Actor-Critic (Lec5) Summary} % Change to "Problem Set X"

% Change the following to your name and UNI.
\newcommand{\name}{Aksel Kretsinger-Walters, adk2164}
\newcommand{\email}{adk2164@columbia.edu}

% NOTE: Defining collaborators is optional; to not list collaborators, comment out the
% line below. Maximum of two collaborators per problem set.
%\newcommand{\collaborators}{Vig Vigerton (\texttt{UNI}), Alice Bob (\texttt{UNI})}
% No collaborators on PS 0

\makeatletter
\def\input@path{{../}{../../}{../../../}} % add as many parents as you need
\makeatother
\input{pset_template_RL.tex} %% DO NOT CHANGE THIS LINE

\begin{document}
\psetheader %% DO NOT CHANGE THIS LINE

\section*{Lecture 5 Summary: Actor--Critic Methods}

\subsection*{Overview}
Actor--critic algorithms combine the strengths of actor-only (policy gradient) and
critic-only (value-based) approaches.
\begin{itemize}
		\item \textbf{Actor-only methods} directly optimize policy parameters $\theta$ using
				Monte Carlo estimates of the policy gradient. They can suffer from high variance and
				lack memory of past estimates.
		\item \textbf{Critic-only methods} (e.g., Q-learning, TD-learning) focus on
				value-function approximation, indirectly deriving a policy but without direct
				gradient optimization guarantees.
\end{itemize}
Actor--critic methods blend both ideas: they estimate value functions (the \emph{critic})
to stabilize and reduce variance in the policy gradient update (the \emph{actor}).

\subsection*{1. Policy Gradient Theorem with Q-Function Approximation}
Recall the policy gradient theorem:
\[
		\nabla_\theta \rho(\pi_\theta)
		= \sum_s d^{\pi_\theta}(s) \, \mathbb{E}_{a \sim \pi_\theta(\cdot|s)}
		\left[ (Q^{\pi_\theta}(s,a) - b^{\pi_\theta}(s)) \nabla_\theta \log \pi_\theta(s,a) \right].
\]
Instead of Monte Carlo estimates of $Q^{\pi_\theta}$, actor--critic methods approximate
the $Q$-function by a parametric model $f_\omega(s,a) \approx Q^{\pi_\theta}(s,a)$.

\paragraph{Theorem (Sutton et al., 1999).}
If $f_\omega$ satisfies the \textbf{compatibility condition}
\[
		\nabla_\omega f_\omega(s,a)
		= \frac{1}{\pi_\theta(s,a)} \nabla_\theta \pi_\theta(s,a)
		= \nabla_\theta \log \pi_\theta(s,a),
\]
and $\omega$ minimizes
\[
		\min_\omega \; \mathbb{E}_{s \sim d^{\pi_\theta}, a \sim \pi_\theta}
		\!\left[ \big(Q^{\pi_\theta}(s,a) - b(s;\theta) - f_\omega(s,a)\big)^2 \right],
\]
then
\[
		\nabla_\theta \rho(\pi_\theta)
		= \mathbb{E}_{s,a \sim \pi_\theta} \!\left[f_\omega(s,a)\,\nabla_\theta \log
		\pi_\theta(s,a)\right].
\]

Thus, $f_\omega$ can safely replace $Q^{\pi_\theta}$ in the policy gradient.

\subsection*{2. Examples of Compatible Function Approximators}
\paragraph{Softmax policy:}
\[
		\pi_\theta(s,a) = \frac{e^{\theta^\top \phi_{s,a}}}{\sum_{a'} e^{\theta^\top \phi_{s,a'}}},
		\quad
		f_\omega(s,a) = \omega^\top (\phi_{s,a} - \sum_{b} \pi_\theta(s,b)\phi_{s,b}).
\]
The critic thus learns a linear function in the same features as the policy—effectively
an approximation of the \emph{advantage function}.

\paragraph{Gaussian policy (continuous actions):}
\[
		\pi_\theta(s,a) = \mathcal{N}(\phi(s)^\top \theta, \sigma^2), \quad
		f_\omega(s,a) = \frac{(a - \theta^\top \phi(s))}{\sigma^2}\phi(s)^\top \omega.
\]

\subsection*{3. Actor--Critic Algorithm}
For each iteration $k$:
\begin{enumerate}
		\item \textbf{Policy evaluation:}
    \[
						\mathbb{E}_{s,a \sim \pi_k}\!\left[(Q^{\pi_k}(s,a) - b_k(s) - f_\omega(s,a))
						\nabla_\omega f_\omega(s,a)\right] = 0.
    \]
		\item \textbf{Policy improvement:}
    \[
						\theta_{k+1} \leftarrow \theta_k + \alpha_k \,
						\mathbb{E}_{s,a \sim \pi_k}\!\left[f_\omega(s,a)\,\nabla_\theta \log \pi_k(s,a)\right].
    \]
\end{enumerate}
This framework updates the critic (value approximation) and actor (policy parameters)
simultaneously.

\subsection*{4. Convergence Guarantees (Sutton et al., 1999)}
Under standard stochastic approximation conditions:
\[
		\sum_k \alpha_k = \infty, \quad \sum_k \alpha_k^2 < \infty,
\]
and bounded second derivatives of $\pi_\theta(s,a)$, we have:
\[
		\lim_{k \to \infty} \nabla_\theta \rho(\pi_{\theta_k}) = 0.
\]
In practice, convergence depends on accurate critic training, stable sampling, and small
gradient noise.

\subsection*{5. Discussion and Modern Extensions}
While theoretically sound, this basic actor--critic scheme is sensitive to approximation
errors. Later works addressed instability and efficiency:
\begin{itemize}
		\item \textbf{Kakade \& Langford (2002)}: Highlighted instability and proposed
				trust-region improvements.
		\item \textbf{TRPO (Schulman et al., 2015a)}: Enforces a KL-divergence constraint on
				policy updates for stability.
		\item \textbf{A3C (Mnih et al., 2016)}: Parallel asynchronous advantage actor--critic.
		\item \textbf{GAE (Schulman et al., 2015b)}: Improves bias--variance tradeoff via
				temporal discounting of advantages.
\end{itemize}

\subsection*{6. Key Insight}
Actor--critic = \textbf{policy gradient + learned baseline}.
The critic provides low-variance, bootstrapped estimates of $Q^\pi$ or $A^\pi$, while the
actor performs gradient ascent using those estimates.

\end{document}
